\documentclass[10pt]{article}
\usepackage[a4paper,margin=20mm]{geometry}
\usepackage{amsmath,amssymb,booktabs,microtype,url,xcolor,float}
\usepackage[section]{placeins}
\usepackage[hidelinks]{hyperref}
\usepackage[numbers,sort&compress]{natbib}
\setlength{\emergencystretch}{2em}
\newcommand{\DevelopmentIncidentCount}{403}
\newcommand{\DevelopmentCausalCount}{256}
\newcommand{\DevelopmentMethodRate}{100.0}
\newcommand{\DevelopmentComparatorRate}{51.6}
\newcommand{\DevelopmentGain}{48.4}
\newcommand{\DevelopmentControlRate}{100.0}
\newcommand{\DevelopmentRollbackRate}{100.0}
\newcommand{\ConfirmationRosterCount}{74}
\newcommand{\SuccessfulSourceCount}{71}
\newcommand{\ValidCausalCount}{226}
\newcommand{\ValidControlCount}{71}
\newcommand{\PrimaryMethodRate}{100.0}
\newcommand{\PrimaryComparatorRate}{49.6}
\newcommand{\PrimaryGain}{50.4}
\newcommand{\PrimaryCILow}{46.6}
\newcommand{\PrimaryCIHigh}{54.6}
\newcommand{\PrimaryP}{0.00001}
\newcommand{\ControlLoss}{0.0}
\newcommand{\ConfirmationStatus}{confirmation gate passed}
\newcommand{\MeanAdapterActions}{7.14}
\newcommand{\MaxAdapterActions}{110}
\newcommand{\MeanRelyAdapterActions}{18.89}


\title{Repair the Boundary, Not the Skills: Contract-Guided Adapter Synthesis for LLM-Agent Skill Composition}
\author{%
\begin{tabular}{c}
Liang Song$^{1,*}$ \qquad Zhai Jiabao$^2$\\[4pt]
\small $^1$ Business School, Xi'an International University, Xi'an 710077, Shaanxi, China\\
\small $^2$ School of Management and Economics, North China University of Water\\
\small Resources and Electric Power, Zhengzhou 450046, Henan, China\\[2pt]
\small $^*$ Corresponding author: \href{mailto:liangsong_1976@126.com}{liangsong\_1976@126.com};\\
\small ORCID: \href{https://orcid.org/0009-0002-6109-0639}{0009-0002-6109-0639}
\end{tabular}%
}
\date{September 2026}

\begin{document}
\maketitle

\begin{abstract}
Reusable Skills let LLM agents preserve successful procedures, but individually
witnessed components can fail when their rely conditions, role bindings, live
state, or validations disagree at composition boundaries. Rewriting a component
widens its regression surface; reverting the whole chain discards unrelated
changes. We study a narrower alternative: compile one executable mediator at the
diagnosed boundary while keeping the component artifacts unchanged. We introduce
Contract-Guided Boundary Adapter Synthesis (CGBAS), which combines typed
rely/guarantee contracts, downstream liveness, and retained successful provenance
to insert a missing witness prefix, specialize a consumer role, or guard a
destructive call before execution. A failed first mechanism revision is retained.
After exposed development on \DevelopmentIncidentCount{} composition incidents,
we froze the method, endpoint, strongest comparator, three model revisions, and
analysis before selecting confirmation tasks. Independent confirmation uses
\ConfirmationRosterCount{} identifier-disjoint ALFWorld and ScienceWorld tasks,
yielding \ValidCausalCount{} state-matched causal conflicts and
\ValidControlCount{} healthy controls. CGBAS achieves
\PrimaryMethodRate\% constrained native repair success versus
\PrimaryComparatorRate\% for the frozen strongest local comparator, a paired
gain of \PrimaryGain{} percentage points (95\% family-stratified task-cluster
bootstrap CI \PrimaryCILow--\PrimaryCIHigh; two-sided task sign-randomization
$p=\PrimaryP$), while healthy-control loss is \ControlLoss{} points. The result
is bounded: CGBAS assumes a correct diagnosis and current successful provenance,
and a single adapter call can contain a long witness prefix. It does not repair
arbitrary production failures or establish that local adaptation always beats
replanning.
\end{abstract}

\noindent\textbf{Keywords:} LLM agents; agent skills; skill composition;
interface repair; design by contract; provenance; ALFWorld; ScienceWorld

\section{Introduction}

LLM-agent Skills package reusable procedures, scripts, tool conventions, and
validation instructions. Reuse turns an isolated successful trajectory into a
software artifact that can be selected and composed in later tasks. Composition,
however, creates failure modes that are not visible when each Skill is evaluated
alone. A consumer may rely on state its predecessor never established; two calls
may bind the same semantic role to different objects; an intermediate call may
destroy state needed later; or a validation may become stale after an effect.

Recent systems jointly predict Skill subsets, counts, and execution order
\citep{zhao2026skillcomposer}, compile Skills into typed execution graphs with
local repair operators \citep{xia2026grasp}, and manage libraries through typed
contracts and compatibility diagnoses \citep{pu2026skillops}. Provenance auditing
attributes reuse to operational evidence \citep{chen2026multitrace}, while
persistent decision histories support repeated Skill evolution
\citep{li2026skillhone}. These advances expose a specific maintenance question
that remains easy to conflate with planning: once a particular interface has
already been diagnosed, must a maintainer rewrite a component or replace the
entire plan, or can the interface itself be repaired?

The distinction matters operationally. A whole-Skill edit invalidates earlier
evidence for all of that Skill's uses. Whole-chain rollback restores a witnessed
version but discards unaffected calls and cannot preserve independent changes.
Rejecting the composition prevents harm but loses capability. A boundary adapter
offers a fourth choice: retain component identities and mediate only the diagnosed
handoff. Such an adapter is useful only if it succeeds in the native environment,
uses supported actions, preserves provenance, and does not damage healthy chains.

We propose Contract-Guided Boundary Adapter Synthesis (CGBAS). Given an ordered
candidate chain, a minimal conflict diagnosis, a public task objective, and a
provenance-linked successful witness chain, CGBAS performs at most one boundary
operation. It can replay the witness prefix omitted before a fixed consumer,
replace a mismatched consumer with its witnessed role-compatible counterpart, or
suppress a diagnosed non-goal destructive call before it fires. Healthy chains
are unchanged. The method therefore treats the adapter as an executable patch to
an interface, not as a newly generated whole plan.

The evaluation is designed around three validity risks. First, native task
success is distinguished from static contract closure. Second, comparison with
LLMs is schema-matched: three revision-pinned open instruction models receive the
same incident, diagnosis, witness catalog, and four allowed operations. Third,
development choices are separated from confirmation. Exposed records from an
earlier diagnosis study are used to revise the mechanism twice; every failed run
is retained. Method, comparator, models, endpoint, gates, environment executors,
and analysis are then hash-locked before a content-blind, zero-overlap task roster
is selected.

This paper contributes:
\begin{itemize}
  \item a boundary-repair formulation with an explicit invariant: preserve all
  unaffected component call identities and introduce at most one adapter call;
  \item CGBAS, which turns rely/guarantee, liveness, role, validation, and
  successful provenance evidence into three executable repair operators;
  \item a causal native evaluation with no repair, a provenance-free pairwise
  heuristic, whole-chain rollback, and three schema-matched local LLM baselines;
  and
  \item an auditable development/confirmation workflow across ALFWorld and
  ScienceWorld, including invalid source tasks, failed revisions, clustered
  inference, action cost, and single-RTX-3090 reproduction materials.
\end{itemize}

We ask: \textbf{RQ1}, can a contract- and provenance-guided boundary adapter
restore native success on independently selected composition conflicts?
\textbf{RQ2}, does it exceed the strongest frozen deployable local comparator
without regressing healthy chains? \textbf{RQ3}, which assumptions and action
costs bound the result?

\begin{figure}[t]
\centering
\fbox{\begin{minipage}{0.94\linewidth}\centering
Exposed development $\rightarrow$ two mechanism revisions $\rightarrow$
method/comparator/analysis hash lock\\[3pt]
$\Downarrow$\\[3pt]
Content-blind zero-overlap roster $\rightarrow$ successful source witnesses
$\rightarrow$ frozen deterministic and LLM programs\\[3pt]
$\Downarrow$\\[3pt]
State-matched candidate/correction execution $\rightarrow$ seven-arm native
replay $\rightarrow$ task-clustered inference
\end{minipage}}
\caption{Evidence order. Confirmation outcomes do not exist until the roster,
source incidents, and every evaluated program are frozen.}
\label{fig:evidence-order}
\end{figure}

\section{Related Work}

\paragraph{Skill selection and composition.}
Generative Skill Composition formalizes selection, count, and order as one
structured prediction problem \citep{zhao2026skillcomposer}. GraSP inserts a
compilation layer between retrieval and execution, representing dependencies as
a typed DAG with node verification and locality-bounded repair
\citep{xia2026grasp}. SkillFuzz uses extracted contracts to prioritize interacting
Skill sets before execution \citep{hu2026skillfuzz}. These systems motivate
composition-level reasoning. Our unit is later and narrower: selection and order
are held fixed, one boundary is already diagnosed, and the outcome is whether a
single mediator can restore a state-matched native task without rewriting the
components.

\paragraph{Skill lifecycle maintenance and auditing.}
SkillOps frames accumulating compatibility and validation defects as Skill
technical debt and supplies a method-agnostic maintenance layer
\citep{pu2026skillops}. SkillHone retains diagnoses, revisions, evidence, and
rejected alternatives across sessions \citep{li2026skillhone}. SkillAudit uses
paired trajectories and fixed structural verification to guide repair without
hidden rewards during evolution \citep{gao2026skillaudit}. SkillCoach separates
process quality from outcome-only success \citep{zhu2026skillcoach}. We adopt
these lifecycle lessons in the experimental design: decision history is retained,
candidate and correction executions are paired, and locality, command support,
provenance, and cost accompany end-task success. The method differs by editing an
execution boundary instead of a Skill document.

\paragraph{Contracts and provenance.}
Preconditions and postconditions make program assumptions explicit
\citep{hoare1969axiomatic,meyer1992contract}; rely/guarantee reasoning distinguishes
what one component assumes from what another establishes in an interfering state
\citep{jones1983rely}. We use these ideas as an operational abstraction over text
simulator calls, not as a proof of arbitrary concurrent programs. SkillTrace
shows that Skill provenance is distributed across expression, implementation,
and operation \citep{chen2026multitrace}. CGBAS uses the operational trace more
narrowly: every added command or role wrapper must cite a retained successful
witness, unless the native environment itself accepts the command.

\section{Problem Formulation}

An ordered Skill composition is $C=(c_1,\ldots,c_n)$. Each call contains immutable
actions $A_i$, an identifier, provenance witnesses, and a contract
\begin{equation}
  K_i=(R_i,G_i,D_i,B_i,V_i),
\end{equation}
where $R_i$ is the rely set, $G_i$ the guaranteed predicates, $D_i$ the clobbered
predicates, $B_i$ semantic role bindings, and $V_i$ validated predicates. Forward
interpretation updates the abstract state as
\begin{equation}
  S_i=(S_{i-1}\setminus D_i)\cup G_i,
\end{equation}
provided $R_i\subseteq S_{i-1}$ and bindings remain compatible. Backward liveness
marks predicates required by any downstream call or by goal completion.

Four diagnosed conflicts are in scope: a \emph{rely gap} when a consumer requires
an absent predicate; a \emph{role mismatch} when an existing role is rebound to a
different object; a \emph{clobber} when an applicable effect destroys a live
predicate; and \emph{stale validation} when such destruction occurs after the
predicate was checked. A healthy chain has none of these conflicts and establishes
goal completion.

Let $W=(w_1,\ldots,w_m)$ be a successful provenance-linked chain for the same
task. An admissible repair $T(C,W)$ may add at most one adapter call at the
diagnosed boundary. Calls not wrapped by that adapter retain their identifiers
and contents. The primary endpoint is constrained native repair success:
\begin{equation}
Y=\mathbb{1}[\text{native goal success}\land\text{locality}\land
\text{provenance}\land\text{adapter-command support}].
\end{equation}
An adapter command is supported if it is copied exactly from retained successful
provenance or accepted by the native environment. Rejected exploratory commands
already present in immutable component calls are reported separately.

\section{Contract-Guided Boundary Adapters}

\subsection{Diagnosis and repair slot}
CGBAS scans the candidate chain once, maintaining abstract state, live bindings,
and validation positions. The first conflict yields its class, evidence, and a
minimal pair of call identifiers. The diagnosis is supplied rather than evaluated
as a contribution in this paper; the intervention starts at the returned slot.

\subsection{Missing-producer prefix}
For a rely gap before consumer $c_j$, CGBAS locates the same immutable consumer
identifier in $W$. It selects the ordered witness calls before that consumer that
are absent from the current prefix, concatenates their actions into one adapter,
and unions their contracts and witness digests. Unlike choosing the shortest call
that establishes one missing predicate, this preserves objective-relevant
processing such as clean, cool, heat, or measure. The operator can therefore be
one call-level modification yet contain many environment actions.

\subsection{Witnessed role wrapper}
For a role mismatch, the live producer bindings are compared with successful
goal-completing witness calls. CGBAS chooses the shortest witnessed consumer whose
non-destination roles match live bindings, replaces the mismatched execution slot
with a wrapper around that call, and records both source and replaced identifiers.
The component artifact is retained; only this invocation is mediated.

\subsection{Pre-execution preservation guard}
Post-effect compensation is impossible when a destructive action terminates an
episode or irreversibly changes the world. For clobber and stale validation,
CGBAS therefore checks effect applicability, goal satisfaction, and downstream
liveness before execution. A diagnosed non-goal destructive call is wrapped by a
zero-action guard. Its original identifier, actions, effect evidence, and witness
remain in adapter metadata. The guard preserves live state rather than attempting
to reconstruct it afterward.

\subsection{Complexity and invariants}
Diagnosis is linear in the calls plus contract-set operations. Witness-prefix
lookup is linear in the witness chain; role selection scans goal-completing witness
calls; a guard requires one forward location scan. The method performs no model
inference. It preserves at most one changed call boundary and leaves healthy
chains unchanged. This is call-level locality, not a constant action-cost claim.

\section{Experimental Method}

\subsection{Development and revision history}
Development uses 403 exposed incidents from a completed earlier diagnostic study:
256 causal conflicts and 147 healthy controls spanning ALFWorld and ScienceWorld
\citep{shridhar2021alfworld,wang2022scienceworld}. Revision 1 inserted the shortest
predicate producer and attempted post-clobber compensation. Its conflict success
was 70.6\%, below the frozen 75\% gate. Failure analysis showed omitted
objective-relevant prefixes and terminal ScienceWorld clobbers. Revision 2 moved
the destructive intervention before execution and retained the whole missing
witness prefix. The first revision's code, outputs, 399 native trajectories, and
four omitted abstentions remain in the artifact.

Before confirmation, we compare Revision 2 with no repair, a pairwise contract
heuristic without provenance or global liveness, whole-chain rollback, and three
local instruction models: Qwen3-4B-Instruct-2507, Phi-4-mini-instruct, and
Mistral-7B-Instruct-v0.3. The models receive identical public inputs and an
identical output schema restricted to no-op, ordered witness-prefix insertion,
witnessed consumer replacement, preservation guard, or abstention. Invalid output
and abstention execute the unchanged candidate and count as failure. A first,
Revision-1-shaped LLM schema was recognized as unfair and retained but excluded;
all three models were rerun under the matched schema before comparator selection.

The strongest deployable local method by development constrained success is
frozen as the primary comparator. Whole-chain rollback is reported as an ordinary
success upper bound, but is not eligible because it violates the primary locality
endpoint. Development advances only with at least 180 causal cases and 60
controls, at least 75\% constrained success, at least a 15-point gain over the
strongest eligible comparator, positive repair gain in every class and both
environments, at most 2-point control loss, and at least 95\% adapter-command
support.

\subsection{Independent confirmation}
The union of task identifiers previously exposed by Papers 1--12 and Paper 13
development is constructed before selection. A salted hash ranks eligible tasks
using only environment, family, identifier, official split, compatibility flag,
and exclusion status. The frozen roster contains 40 untouched ALFWorld train
instances balanced across five families and all 34 eligible untouched
ScienceWorld use-thermometer test variations. These ALFWorld tasks support
task-level transfer under zero overlap, not a claim about the official test split.

Native expert or gold trajectories are collected after the roster freeze. Failed
sources are retained and excluded, not replaced. For each successful source, the
constructor creates a healthy control and four possible one-boundary mutations:
omit the producer prefix, substitute a same-family donor consumer with a different
subject, insert an applicable non-goal clobber, or validate immediately before the
clobber. Mutation specifications and public incidents are frozen before candidate
or correction execution.

All deterministic and LLM program files are then frozen. Only afterward do we
execute the candidate and matched correction from the same native initial state.
A causal label requires successful source provenance, candidate failure,
matched-correction success, two matching initial states, and no technical error.
A control requires source and unchanged-candidate success with a matching state.
Invalid mutations are excluded rather than relabeled as controls. Every program
arm is independently replayed on every planned event; technical errors remain
failures.
Figure~\ref{fig:evidence-order} summarizes this ordering; file hashes make each
boundary between stages auditable.

\subsection{Inference and gates}
The primary comparison is paired constrained success between CGBAS and the frozen
strongest comparator over valid causal conflicts. Because multiple conflict
classes derive from one source task, inference clusters them by task. We use
10,000 family/environment-stratified task bootstrap replicates for a 95\% interval
and a two-sided 100,000-replicate paired task sign-randomization test
\citep{efron1994bootstrap}. Confirmation requires at least 60 source tasks, 180
causal conflicts, 60 controls, and 30 examples per conflict class; CGBAS success
at least 75\%; at least a 15-point comparator gain with the interval above zero
and $p\leq0.01$; positive repair gain over no repair in every class and both
environments; at most 2-point control loss; and complete locality, provenance,
initial-state, command-support, and technical-error audits.

\section{Results}

\subsection{Development gate}
\begin{table}[t]
\centering
\small
\caption{Exposed development results. The strongest eligible local comparator is frozen before confirmation.}
\begin{tabular}{lrrr}
\toprule
Arm & Native (\%) & Constrained (\%) & Control (\%) \\
\midrule
CGBAS & 100.0 & 100.0 & 100.0 \\
Mistral-7B & 51.6 & 51.6 & 100.0 \\
Phi-4-mini & 46.5 & 46.5 & 100.0 \\
Qwen3-4B & 35.2 & 35.2 & 100.0 \\
Pairwise contract & 10.2 & 9.8 & 100.0 \\
No repair & 0.0 & 0.0 & 100.0 \\
Whole-chain rollback & 100.0 & 0.0 & 100.0 \\
\bottomrule
\end{tabular}
\end{table}


Revision 2 reaches \DevelopmentMethodRate\% constrained success across
\DevelopmentCausalCount{} causal conflicts. Mistral-7B is selected automatically
as the strongest eligible comparator at \DevelopmentComparatorRate\%, leaving a
\DevelopmentGain{}-point development difference. Healthy-control success remains
\DevelopmentControlRate\%. Whole-chain rollback reaches
\DevelopmentRollbackRate\% ordinary success but fails the locality endpoint.

\subsection{Independent primary result}
\begin{table}[t]
\centering
\small
\caption{Independent confirmation. Invalid mutations are excluded through one common state-matched label filter.}
\begin{tabular}{lrrrr}
\toprule
Arm & $n$ & Native (\%) & Constrained (\%) & Control (\%) \\
\midrule
CGBAS & 226 & 100.0 & 100.0 & 100.0 \\
Mistral-7B & 226 & 49.6 & 49.6 & 100.0 \\
Phi-4-mini & 226 & 42.9 & 42.9 & 100.0 \\
Qwen3-4B & 226 & 36.7 & 36.7 & 100.0 \\
Pairwise contract & 226 & 8.4 & 7.5 & 100.0 \\
No repair & 226 & 0.0 & 0.0 & 100.0 \\
Whole-chain rollback & 226 & 100.0 & 0.0 & 100.0 \\
\bottomrule
\end{tabular}
\end{table}


Of \ConfirmationRosterCount{} roster tasks, \SuccessfulSourceCount{} yield valid
source trajectories. The state-matched validity rules retain
\ValidCausalCount{} causal conflicts and \ValidControlCount{} controls. CGBAS
achieves \PrimaryMethodRate\% constrained success, compared with
\PrimaryComparatorRate\% for the frozen comparator. The paired difference is
\PrimaryGain{} points; its 95\% clustered interval is
\PrimaryCILow--\PrimaryCIHigh{} and the paired task sign-randomization test gives
$p=\PrimaryP$. The predeclared confirmation status is \ConfirmationStatus{}.

\subsection{Conflict classes, environments, and cost}
\begin{table}[t]
\centering
\small
\caption{Constrained native success by conflict class.}
\begin{tabular}{lrrrr}
\toprule
Conflict & $n$ & CGBAS (\%) & Frozen comparator (\%) & No repair (\%) \\
\midrule
rely gap & 71 & 100.0 & 43.7 & 0.0 \\
role mismatch & 71 & 100.0 & 0.0 & 0.0 \\
clobber & 42 & 100.0 & 100.0 & 0.0 \\
stale validation & 42 & 100.0 & 92.9 & 0.0 \\
\bottomrule
\end{tabular}
\end{table}

\begin{table}[t]
\centering
\small
\caption{Causal constrained success by native environment.}
\begin{tabular}{lrrrr}
\toprule
Environment & $n$ & CGBAS (\%) & Frozen comparator (\%) & No repair (\%) \\
\midrule
alfworld & 128 & 100.0 & 64.1 & 0.0 \\
scienceworld & 98 & 100.0 & 30.6 & 0.0 \\
\bottomrule
\end{tabular}
\end{table}


Class-level repair gains over no repair are reported separately from the overall
comparison with the strongest baseline. This prevents a baseline tie in one
mechanically easier class from being hidden by averaging. Environment results
likewise preserve all conflict instances while clustering inference at the source
task. CGBAS proposes \MeanAdapterActions{} adapter actions per causal repair on
average and at most \MaxAdapterActions{}. The rely-gap cost is
\MeanRelyAdapterActions{} actions on average, making explicit the difference
between one call-level boundary edit and a short execution.

\subsection{Integrity audits}
\begin{table}[t]
\centering
\small
\caption{Predeclared confirmation gates and integrity checks.}
\begin{tabular}{lr}
\toprule
Frozen check & Result \\
\midrule
minimum tasks & Pass \\
minimum causal conflicts & Pass \\
minimum controls & Pass \\
minimum per conflict class & Pass \\
minimum constrained success & Pass \\
minimum gain over frozen comparator & Pass \\
bootstrap interval above zero & Pass \\
paired randomization p at most alpha & Pass \\
positive repair gain every conflict class & Pass \\
positive repair gain both environments & Pass \\
control loss at most two points & Pass \\
minimum 95 percent adapter command support & Pass \\
all primary repairs local & Pass \\
all primary provenance complete & Pass \\
all primary initial states match & Pass \\
no primary technical errors & Pass \\
zero overlap roster & Pass \\
\bottomrule
\end{tabular}
\end{table}


The roster has zero task-identifier overlap with the prior-exposure union. All
program artifacts precede candidate and adapter outcomes. Model revision
identifiers, prompts, method code, comparator, executors, endpoint, gates, and
analysis hashes match the pre-confirmation lock. Invalid sources and mutations remain visible in the released
label table, and all planned program outcomes remain in arm denominators until the
common validity filter is applied.

\section{Discussion}

The main practical result is conditional. When a maintainer has a reliable
boundary diagnosis and retained successful provenance, the repair decision need
not become a whole-Skill rewrite. Rely gaps require restoring the causal prefix,
not merely satisfying one abstract predicate. Destructive calls must be guarded
before they fire when compensation may be impossible. Role failures can be
repaired by a witnessed consumer invocation rather than free-form command
generation. These mechanisms explain why contract-only adjacency and local LLM
selection can underperform despite receiving the same diagnosed slot.

Whole-chain rollback remains important. It can recover ordinary task success and
may be preferable when provenance is cheap to replay, independent changes need not
be preserved, or the diagnosis is uncertain. CGBAS contributes a different point
in the engineering tradeoff: it preserves unaffected call identities and exposes
exactly what was inserted or suppressed. Call locality should not be confused
with runtime economy, because a missing producer can require a long witness prefix.

The result also suggests a lifecycle architecture. Composition diagnostics can
emit a typed boundary and evidence; a provenance store can supply compatible
witness slices; an adapter compiler can produce a versioned invocation wrapper;
and release gates can replay both conflicts and healthy controls. The immutable
component and explicit adapter make rollback, audit, and later garbage collection
simpler than silently editing a shared Skill.

\section{Threats to Validity}

\paragraph{Construct validity.}
The four mutations are controlled interface regressions generated from successful
traces. They do not estimate the prevalence of those causes in production. Causal
labels require native candidate failure and matched correction success, but the
same mutation operators were used in development and confirmation. The evidence
therefore supports task transfer, not operator discovery.

\paragraph{Information advantage.}
CGBAS receives a correct diagnosis and a successful witness catalog. This is
appropriate for the retained-provenance maintenance setting but stronger than the
information available after an arbitrary failure. The LLM comparators receive the
same catalog, diagnosis, and operation set; they are not general autonomous agents.

\paragraph{External validity.}
ALFWorld and ScienceWorld provide deterministic text interfaces, native goals,
and replayable resets. Production APIs may drift, exhibit side effects, or lack
exact reset and outcome oracles. ScienceWorld confirmation covers one eligible
family after strict zero-overlap exclusion. Results should be extended to tools,
web tasks, and longer asynchronous compositions before broader claims.

\paragraph{Statistical validity.}
Derived conflict instances are not independent; all inference clusters by source
task and stratifies by family/environment. Nevertheless, 74 roster tasks are
modest, and class-specific estimates are less precise than the primary aggregate.
Thresholds and analysis were frozen before confirmation, and no failed source was
replaced.

\paragraph{Implementation validity.}
Contracts are compiled from successful traces, and native command support is
audited only for adapter-added actions. Original Skills may contain exploratory
commands that are absent from an environment's returned action list. We report
such behavior but do not attribute it to the adapter. Hash locks, append-only
traces, initial-state signatures, and a standalone replay bundle reduce but do not
eliminate implementation risk.

\section{Conclusion}

CGBAS repairs an already diagnosed LLM-agent Skill interface by compiling one
provenance-backed boundary adapter. It restores omitted witness prefixes,
specializes consumer roles, or prevents live-state destruction before execution,
while retaining unaffected component identities. A frozen, zero-overlap native
evaluation tests whether that mechanism exceeds no repair, contract-only local
repair, whole-chain rollback under a locality constraint, and three schema-matched
LLM selectors. The conclusion is intentionally scoped to controlled interface
regressions with current successful provenance. Within that scope, repairing the
boundary is an auditable alternative to rewriting reusable Skills; outside it,
diagnosis uncertainty, provenance drift, irreversible effects, and long replay
cost remain open problems.
\section*{Author Contributions}
Liang Song: conceptualization, methodology, software, investigation, formal
analysis, visualization, and writing--original draft. Zhai Jiabao: validation,
data curation, reproducibility review, and writing--review and editing. Both
authors approved the submitted version and accept responsibility for the work.

\section*{Funding}
No funding was received for this research.

\section*{Declaration of Competing Interest}
The authors declare that they have no known competing financial or non-financial
interests that could have influenced the work reported in this paper.

\section*{Data and Code Availability}
The study-specific mutation records, paired traces, analysis scripts, and
reproduction materials are prepared for release in a public repository with a
DOI. The repository DOI and permanent URL will be inserted before submission.
ALFWorld, ScienceWorld, and model-weight assets remain subject to their
respective upstream licenses and are identified rather than redistributed.

\section*{Acknowledgements}
None.

\section*{Generative AI Disclosure}
No generative AI system was used to generate experimental data, labels, or
results. Any language assistance used during manuscript preparation will be
disclosed in accordance with the journal's policy.

\bibliographystyle{plainnat}
\bibliography{references}
\end{document}
